\SourceLeaf{417}%
arc $aa'$ and let $\lambda,\Lambda$ denote in relation to that point what $\alpha,A$ and $\alpha',A'$ denote in relation to $a$ and $a'$. If now $\alpha\lambda$ or $\Lambda-A$ is proportional to the part of the arc $al$, we shall say that $aa'$ is uniformly curved along its entire length, and call
\[
\frac{\Lambda-A}{al}
\]
the measure of curvature or simply the curvature. It is easy to see that this occurs only when $aa'$ really is a circular arc and that by our definition its curvature then comes out $=\pm\dfrac1r$ if $r$ denotes the radius. Insofar as $r$ is always taken as positive, the upper or lower sign applies according as the centre lies to the right or left of the arc $aa'$ ($a$ being taken as the initial point, $a'$ as the endpoint, and the directions in the auxiliary circle being counted from left to right). Changing one of these conditions changes the sign, and changing two restores it.

If, however, $\Lambda-A$ is not proportional to $al$, we call the arc nonuniformly curved, and the quotient
\[
\frac{\Lambda-A}{al}
\]
may then be called its mean curvature. Curvature without qualification, however, always presupposes specifying the location and is defined by the mean curvature of an element there. It is therefore
\[
=\frac{d\Lambda}{dal}.
\]
We thus see that arc, Amplitudo and curvature stand in the same relation to one another as time, motion and velocity, or volume, mass and density. The reciprocal of the curvature, namely
\[
\frac{dal}{d\Lambda}
\]
is called the radius of curvature at $l$. Under the assumptions above the curve at this point is called concave towards the right and convex towards the left if the value of its curvature or radius of curvature is positive. If it is negative, the reverse holds.
